%% CVFEM HEX8 Navier-Stokes: the operator, component by component, and the micro-kernels.
%%
%% This is the LONG-FORM companion to wip/paper. The paper is hard-limited to ten pages and
%% carries the performance argument; it cannot carry the discretisation. This document does the
%% opposite: it states every component of the operator in closed form, says which micro-kernel
%% evaluates it for each scheme and layout, and names the file and function that implements it so
%% a reader can go from a formula to the code that runs.
%%
%% TWO RULES THIS DOCUMENT KEEPS.
%%
%% Every formula is the one the code evaluates, not a textbook form it resembles. Where the
%% implementation departs from the usual presentation -- the Rhie-Chow time scale, the
%% Darwish-Moukalled algebra, the limiter subgradients -- the departure is stated and the reason
%% given, because those are exactly the places where reading the paper's equations and the source
%% side by side would otherwise mislead.
%%
%% No ABSOLUTE timing is typed in here. Throughputs belong to the paper, which generates them from
%% the measurement files. What this document carries are (a) COUNTS -- loads per surface, live
%% values, lines of generated kernel -- which are properties of the written kernel, and (b) RATIOS
%% between alternatives measured against each other in ONE job, which are the evidence for the
%% selection table of Section 6 and cannot be stated without them. Every ratio names its job id,
%% because the tree moved under these measurements and a ratio from one state is not comparable
%% with an absolute from another.
\documentclass[11pt,a4paper]{article}

\usepackage[margin=1in]{geometry}
\usepackage{amsmath,amssymb}
\usepackage{booktabs}
\usepackage{microtype}
\usepackage[hidelinks]{hyperref}
\usepackage{xcolor}
\usepackage{listings}

\lstset{basicstyle=\ttfamily\footnotesize,columns=fullflexible,keepspaces=true,
        breaklines=true,frame=single,framesep=4pt,rulecolor=\color{black!25},
        showstringspaces=false,language=C++,commentstyle=\color{black!55}\itshape}

% The operator's recurring symbols, defined once so the document cannot drift from itself.
\newcommand{\scs}{\mathrm{s}}                      % a sub-control surface
\newcommand{\Avec}{\mathbf{A}}                     % its area vector
\newcommand{\dvec}{\mathbf{d}}                     % its node-to-node edge vector
\newcommand{\uvec}{\mathbf{u}}
\newcommand{\xvec}{\mathbf{x}}
\newcommand{\Jmat}{\mathsf{J}}                     % the element Jacobian
\newcommand{\adj}{\operatorname{adj}}
\newcommand{\mdot}{\dot{m}}
\newcommand{\code}[1]{\texttt{\small #1}}

\title{The CVFEM HEX8 Navier--Stokes operator\\
       \large Components, closed forms, and the micro-kernel for each scheme}
\author{}
\date{}

\begin{document}
\maketitle

\begin{abstract}
\noindent
This is the long-form companion to the performance paper. It states each component of the
control-volume finite-element HEX8 Navier--Stokes operator in the form the code evaluates, and for
each convective scheme it describes the micro-kernel that evaluates it: the loop structure, what it
reads per sub-control surface, how it is vectorised, and which of the two mesh layouts changes what.
Every section names the file and function it describes. Where the implementation departs from the
textbook statement of a scheme --- and it does, in the Rhie--Chow time scale, in the
Darwish--Moukalled algebra, and in every limiter derivative --- the departure and its reason are
given, since those are the places where the source and a naive reading of the equations part
company.
\end{abstract}

\tableofcontents

\section{The discretisation}
\label{sec:disc}

The operator is a control-volume finite-element (CVFEM) discretisation of incompressible
Navier--Stokes on trilinear hexahedra. Each of the eight nodes of an element owns one
sub-control volume, and the element contributes to the flux balance of those volumes across
twelve \emph{sub-control surfaces} (SCS). The unknowns are collocated: each node carries
$(u_x,u_y,u_z,p)$, so an element holds $8\times 4 = 32$ degrees of freedom.

\subsection{Sub-control surfaces}

The twelve surfaces are enumerated in \code{CVFEM\_HEX8\_SCS}
(\code{src/upwind/cvfem\_hex8\_ns\_upwind\_kernels.hpp}). Each entry is a triple
$(i, j, \Avec^{\mathrm{ref}})$: the two nodes whose sub-control volumes the surface separates,
and the surface's area vector on the reference element. The enumeration is grouped
\emph{four per reference direction}:

\begin{center}
\begin{tabular}{@{}llll@{}}
\toprule
direction & $\scs$ & node pairs $(i,j)$ & $\Avec^{\mathrm{ref}}$ \\
\midrule
$\xi$   & $0,1,2,3$    & $(0,1)\ (3,2)\ (4,5)\ (7,6)$ & $(\tfrac14,0,0)$ \\
$\eta$  & $4,5,6,7$    & $(0,3)\ (1,2)\ (4,7)\ (5,6)$ & $(0,\tfrac14,0)$ \\
$\zeta$ & $8,9,10,11$  & $(0,4)\ (1,5)\ (2,6)\ (3,7)$ & $(0,0,\tfrac14)$ \\
\bottomrule
\end{tabular}
\end{center}

Two facts follow from this table and are used throughout the implementation. First,
$\scs \,/\, 4$ is the reference direction, and because $\scs$ is a template parameter in every
face kernel the direction is a compile-time constant. Second, each pair $(i,j)$ is
\emph{edge-adjacent}: the two nodes differ in exactly one logical coordinate, the one naming the
group. Section~\ref{sec:geom} shows that this is what makes the element Jacobian sufficient to
describe the surface geometry on the affine path.

Every node appears in exactly three of the twelve surfaces, once per direction group. That
multiplicity is the reason the element's four accumulators per node ($r_x, r_y, r_z$ and the
continuity row) are live across the whole surface loop, which
Section~\ref{sec:mk:structure} returns to as the binding constraint on the kernel.

\subsection{The residual}

Writing $R_a$ for the contribution to node $a$'s sub-control volume, the element residual is a
sum over the twelve surfaces of a flux times an area, plus the volume terms:
\begin{equation}
R_a \;=\; \sum_{\scs \,:\, i(\scs)=a} \mathbf{F}_\scs
      \;-\; \sum_{\scs \,:\, j(\scs)=a} \mathbf{F}_\scs
      \;+\; \text{(transient)}_a \;+\; \text{(boundary)}_a .
\label{eq:residual}
\end{equation}
The antisymmetry is exact and structural: a surface adds its flux to node $i$ and subtracts the
same value from node $j$, so the element conserves by construction regardless of what
$\mathbf{F}_\scs$ is. Every kernel in this document preserves it literally --- the flux is
computed once per surface and applied with opposite signs --- which is why no scheme described
here can break conservation, only accuracy.

The surface flux has three parts, treated in Section~\ref{sec:components}: a viscous part, a
convective part carrying the mass flux, and the pressure--velocity coupling that fixes the mass
flux itself. The convective part is where the schemes differ, and Section~\ref{sec:schemes}
takes those one at a time.

\section{Geometry: the affine and isoparametric variants}
\label{sec:geom}

The reference element is the unit cube $[0,1]^3$ with trilinear shape functions. The Jacobian
$\Jmat(\xi) = \partial\xvec/\partial\xi$ is in general a function of position; the two variants
differ in whether they admit that.

\subsection{What each variant commits to}

The \emph{isoparametric} variant evaluates the geometry where it is needed, so $\Jmat$, its
adjugate and the surface areas all vary within the element
(\code{cvfem\_hex8\_geom\_at}, and the \code{\_isoparam} kernels).

The \emph{affine} variant evaluates $\Jmat$ once, at the reference centre
$\xi = (\tfrac12,\tfrac12,\tfrac12)$, and uses it everywhere
(\code{cvfem\_hex8\_affine\_adj}). For a parallelepiped this is exact. For a general trilinear
hexahedron it is an approximation, and it is the \emph{defining} approximation of the variant:
having made it, the variant's geometry \emph{is} that one constant Jacobian, and a kernel on this
path has no entitlement to any other geometric information. This is a stronger statement than it
first appears, and Section~\ref{sec:geom:edges} shows where ignoring it introduced a real bug.

What is stored per element is the adjugate and the determinant, not $\Jmat$ itself. With
$\adj(\mathsf{A}) = \det(\mathsf{A})\,\mathsf{A}^{-1}$ the code keeps
\code{cof0..cof8} $=\adj(\Jmat)$ row-major and \code{det} $=\det\Jmat$.

\subsection{Surface areas}

The physical area vector of surface $\scs$ in direction $q = \scs/4$ is the reference area vector
pushed forward by the cofactor matrix. Because $\Avec^{\mathrm{ref}}$ is $\tfrac14$ along axis
$q$, this reduces to a scaling of one row of the adjugate:
\begin{equation}
\Avec_\scs \;=\; \tfrac14 \bigl(\adj\Jmat\bigr)_{q,\cdot} ,
\end{equation}
which is exactly the nine lines \code{Ax0 = qtr*c0} \dots{} \code{Az2 = qtr*c8} in the geometry
loop of each sweep. All four surfaces of a direction group share one area vector, so the sweep
holds three of them, not twelve.

\subsection{Edge vectors, and why the Jacobian suffices}
\label{sec:geom:edges}

Several components need $\dvec_\scs$, the vector from node $i$ to node $j$ across the surface.
The obvious implementation differences the node coordinates, $\dvec_\scs = \xvec_j - \xvec_i$,
and that is what the code did until recently.

On the affine path it is the wrong choice, and not because it is inaccurate --- because it is
\emph{inconsistent}. Evaluating the trilinear Jacobian at the reference centre gives, for
direction $q$,
\begin{equation}
\bigl(\Jmat\bigr)_{\cdot,q} \;=\; \frac{1}{4}\!\!\sum_{\scs \,\in\, \text{group } q}\!\!
      \bigl(\xvec_{j(\scs)} - \xvec_{i(\scs)}\bigr),
\label{eq:col-is-mean}
\end{equation}
that is, column $q$ of the centre Jacobian is exactly the arithmetic mean of the four true edge
vectors of group $q$. This holds for \emph{any} trilinear hexahedron, not only a parallelepiped:
it is an identity of the trilinear basis evaluated at the centre, and it was verified
symbolically against the code's own adjugate expressions on cubes, sheared parallelepipeds and
randomly distorted elements, agreeing to $10^{-16}$ in every case.

So the affine variant has already committed to that mean when it computes the flux. A kernel that
then differences node coordinates for $\dvec$ is using a \emph{second, different} geometry for one
term while the rest of the same kernel uses the first. On a cube the two coincide exactly; on a
distorted element the true edges differ from their mean by $10\%$ to over $100\%$, which is a
measure of how far outside its validity the affine variant has been pushed, not of an error in
either formula.

The affine path therefore takes
\begin{equation}
\dvec_\scs \;=\; \bigl(\Jmat\bigr)_{\cdot,\,\scs/4},
\end{equation}
and recovers $\Jmat$ from what is already in registers using the $3\times 3$ identity
\begin{equation}
\adj\bigl(\adj \mathsf{A}\bigr) \;=\; \det(\mathsf{A})\,\mathsf{A}
\qquad\Longrightarrow\qquad
\Jmat \;=\; \frac{\adj(\adj\Jmat)}{\det\Jmat} ,
\label{eq:adj-adj}
\end{equation}
implemented in \code{cvfem\_hex8\_affine\_edge\_cols}. The cofactors and $1/\det$ are already live
in the geometry loop that builds the area vectors, so the three columns cost roughly forty flops
per element and no loads, and the four surfaces of a group share one column. The isoparametric
kernels keep true per-surface coordinates, which is correct where the geometry genuinely varies.

\paragraph{A consistency obligation this creates.} Every \emph{producer} of precomputed data must
use the same $\dvec$ as every \emph{consumer}. The coefficient table
(\code{cvfem\_hex8\_build\_rc\_coeff}) and the partially assembled tangent
(\code{cvfem\_hex8\_build\_pa\_tangent}) are both producers, and both had to move to
Eq.~\eqref{eq:adj-adj} with the kernels. Missing the second of them put the stored tangent and the
kernel reading it on different geometries, which \code{cvfem\_pa\_tangent\_test} caught at
$1.4\times10^{-5}$ --- five orders of magnitude above round-off, and invisible to every
cube-based check, because on a cube Eq.~\eqref{eq:col-is-mean} makes the two definitions
identical. Any future change here needs a non-affine oracle.

\section{The operator's components}
\label{sec:components}

Each component below is given in the form the code evaluates, with the function that evaluates it.
Throughout, $\scs$ is a sub-control surface with nodes $i,j$, area vector $\Avec$, edge vector
$\dvec$ (Section~\ref{sec:geom:edges}), and $\bar{\uvec} = \tfrac12(\uvec_i + \uvec_j)$ is the
face-average velocity.

\subsection{The viscous term}

The viscous flux uses the element's velocity gradient, obtained from the nodal values by the
reference-gradient-then-pushforward route shared with every other gradient in the code
(\code{cvfem\_hex8\_grad\_at}): the reference gradient is contracted from the nodal values, then
mapped by $\adj\Jmat/\det\Jmat$. In the sweeps the nine components are held in
\code{g00v}\dots\code{g22v}, one lane-major array each, built once per element in the geometry
loop. The surface then contributes $\mu\,\nabla\uvec\cdot\Avec$ per velocity component, which is
\code{cvfem\_hex8\_visc\_dir\_simd} instantiated once per direction group --- three calls, each
reusing one area vector across its four surfaces.

The viscous term is scheme-independent and layout-independent. It is also the reason the Jacobian
is not bit-identical to the residual under reassociation: the viscous contribution is already in
the accumulator when the surface loop starts, so summing a node's three surface contributions in a
register before adding changes the order of additions. That is round-off, but it is why the
fingerprint moves when the surface loop is restructured.

\subsection{The convective flux and the mass flux}

The surface mass flux is
\begin{equation}
\mdot_\scs \;=\; \rho\,\bar{\uvec}\cdot\Avec \;-\; \underbrace{G_\scs\,c_\scs}_{\text{Rhie--Chow}},
\label{eq:mdot}
\end{equation}
with the second term deferred to Section~\ref{sec:rc}. The convected momentum uses an upwind
split of $\mdot$:
\begin{equation}
\mdot^{+} = \tfrac12\bigl(\mdot + |\mdot|_\epsilon\bigr), \qquad
\mdot^{-} = \tfrac12\bigl(\mdot - |\mdot|_\epsilon\bigr),
\end{equation}
so that $\mdot^{+}$ vanishes for inflow and $\mdot^{-}$ for outflow, and the momentum flux is
\begin{equation}
\mathbf{F}^{\text{conv}}_\scs \;=\; \mdot^{+}\uvec_i \;+\; \mdot^{-}\uvec_j
      \;+\; \tfrac12\bigl(p_i + p_j\bigr)\Avec \;+\; \mathbf{h}_\scs ,
\label{eq:flux}
\end{equation}
where the last term is the higher-order correction of Section~\ref{sec:schemes} and is zero for
first-order upwinding. The continuity row of the surface is $\mdot_\scs$ itself, carried with the
same antisymmetry as the momentum rows.

\paragraph{The Harten band.} $|\mdot|_\epsilon$ is the absolute value with an optional smoothing
band of half-width $\epsilon$ (\code{cvfem\_upwind\_abs}). With $\epsilon = 0$ it is the hard
switch. The band exists because the hard switch is non-differentiable exactly where $\mdot$
changes sign, which a Newton method visits. It is a \emph{compile-time} flag in every kernel
(\code{EPS}), not a runtime argument: a runtime zero would leave the comparison in the vector body,
and a guard of that shape in this loop was measured at $1.83\times$.

\subsection{Rhie--Chow pressure--velocity coupling}
\label{sec:rc}

Collocated storage admits a checkerboard pressure mode, because the face mass flux built from
$\bar{\uvec}$ alone does not see pressure oscillations at the mesh scale. Rhie--Chow adds to
$\mdot$ a term proportional to the difference between the true pressure difference across the
surface and the one implied by the nodal pressure gradients:
\begin{equation}
c_\scs \;=\; \bigl(p_j - p_i\bigr) \;-\; \tfrac12\bigl(\nabla p_i + \nabla p_j\bigr)\cdot\dvec ,
\label{eq:corr}
\end{equation}
which vanishes identically for a pressure field whose nodal gradient is exact, and is $O(h^2)$
for a smooth one --- so the term is a filter on the mesh-scale mode and not a modification of the
continuous equations. The nodal pressure gradient is reconstructed by its own pass
(Section~\ref{sec:recon}).

The coefficient $G_\scs$ is $V/a_P$ in the usual presentation --- a cell volume over a momentum
diagonal. The implementation (\code{cvfem\_hex8\_rhie\_chow\_mdot\_coeff}) evaluates
\begin{equation}
G_\scs \;=\; \alpha\,\frac{|\Avec|^2}{\Avec\cdot\dvec}\,
      \Bigl[\bigl(2a_0/\Delta t\bigr)^2 + \frac{4|\bar{\uvec}|^2}{h^2}
            + \Bigl(\frac{4\nu}{h^2}\Bigr)^{\!2}\Bigr]^{-1/2},
\qquad h^2 = |\dvec|^2 ,
\label{eq:coeff}
\end{equation}
with $\alpha$ the user scale. The bracket is the Shakib--Tezduyar combined time scale, and the
choice matters: the earlier implementation used the diffusion limit alone,
$\alpha h^2/(2\mu)$, which is $V/a_P$ computed as though the momentum diagonal were purely
viscous. That overstates $G$ by roughly the cell P\'eclet number $\rho|u|h/\mu$ --- about $31\times$
on a lid-driven cavity at $Re=1000$, about $159\times$ on a backward-facing step at
$Re_h = 5100$ --- at which point the "small" correction the derivation assumes is setting a
substantial part of the face mass flux by pressure smoothing rather than by the flow. Note that
only $|\bar{\uvec}|^2$ is needed, never $|\bar{\uvec}|$, so no square root is spent on the
velocity.

\paragraph{What is tabulated and what is not.} $G_\scs$ depends on the state through
$|\bar{\uvec}|^2$, so it is rebuilt once per Newton step, not once per matvec, and stored
element-major as \code{rc\_coeff}$[e\cdot 12 + \scs]$. The velocity sensitivity needed by the
Jacobian,
\begin{equation}
w_\scs \;=\; \frac{4\,u^2_{\text{scale}}\,(\Avec\cdot\dvec)^2}{h^2\,\alpha^2\,|\Avec|^4},
\label{eq:wdu}
\end{equation}
is by contrast \emph{purely geometric} plus two uniforms --- no state at all --- and is tabulated
beside the coefficient as \code{rc\_w}. It was previously recomputed in the surface loop, three dot
products and a division per surface, for a quantity that never varies with the solution.

\subsection{Transient and boundary terms}

The transient term is a volume contribution per sub-control volume, applied as a separate pass
(\code{apply\_transient\_action\_pass}) so that a steady solve pays nothing for it. The boundary
closure is likewise a separate pass over boundary sub-control surfaces. Both are
scheme-independent; the paper's completeness ladder reports the operator with each switched on in
turn, and any verification of the Jacobian must include whichever passes the timed apply included
--- omitting them makes a finite-difference check report the closure as a staging error.

\section{The convective schemes}
\label{sec:schemes}

The schemes differ only in $\mathbf{h}_\scs$, the higher-order correction in
Eq.~\eqref{eq:flux}. All five share the mass flux, the upwind split and the pressure term, which
is why the Jacobian stays first order unless the correction is differentiated
(Section~\ref{sec:jac}).

\subsection{Deferred correction}

The correction reconstructs the face value from the donor node and its nodal velocity gradient,
and takes only the \emph{difference} from the first-order upwind value:
\begin{equation}
\mathbf{h}_\scs \;=\; \mdot^{+}\,\delta_i \;+\; \mdot^{-}\,\delta_j ,
\qquad
\delta_a \;=\; \Lambda\!\left(u_a,\; u_{a'},\; \nabla u_a \cdot (\xvec_\scs^{c} - \xvec_a)\right),
\label{eq:defcor}
\end{equation}
where $\xvec_\scs^{c}$ is the surface centroid, $a'$ is the surface's other node, and $\Lambda$ is
the limiter. With $\Lambda$ the identity this is unlimited second-order upwinding; the three
limiters below are the three non-trivial choices of $\Lambda$. Writing it as a difference from
first order is what makes it \emph{deferred}: the correction enters the residual and is absent
from the lagged Jacobian, so the linear system keeps the first-order sparsity and the
higher-order accuracy is recovered through the Newton iteration.

The unlimited increment $\Delta = \nabla u_a \cdot (\xvec^c_\scs - \xvec_a)$ is what each limiter
acts on, and $[\,\ell, h\,] = [\min(u_a,u_{a'}),\,\max(u_a,u_{a'})]$ is the interval the
surface's own two nodes span --- a two-node stencil, not a full one-ring.

\subsection{Scheme 0: unlimited}

$\Lambda(\Delta) = \Delta$. The arm every other is measured against, and the solver's default.
Smooth, differentiable, and does not preserve the bound: a reconstruction can place the face value
outside the interval its nodes span.

\subsection{Scheme 1: bounded-face clip}

Barth--Jespersen's bound on the two-node stencil
(\code{cvfem\_limiter\_clip\_inc}):
\begin{equation}
\Lambda(\Delta) \;=\; \operatorname{clip}\bigl(u_a + \Delta,\; \ell,\; h\bigr) - u_a .
\end{equation}
The face value cannot introduce a new extremum, so the bound holds by construction. It is kept as
the reference arm rather than recommended: it clips at a smooth extremum too, and its active set
flips --- an ulp of change in $\Delta$ can move the result by a finite amount, which is exactly
what makes its derivative delicate (Section~\ref{sec:jac:subgrad}).

\subsection{Scheme 2: Venkatakrishnan}

A smooth limiter (\code{cvfem\_venkata\_inc}). With $D^{+}$ the signed room to the bound in the
direction of the increment,
\begin{equation}
D^{+} = \begin{cases} h - u_a, & \Delta \ge 0\\ \ell - u_a, & \Delta < 0\end{cases}
\qquad
\Lambda(\Delta) \;=\; \frac{(D^{+})^2 + 2\Delta D^{+} + \varepsilon^2}
                          {(D^{+})^2 + 2\Delta^2 + \Delta D^{+} + \varepsilon^2}\;\Delta .
\label{eq:venkat}
\end{equation}
The denominator is a sum of squares, $(D^{+} + \Delta/2)^2 + \tfrac74\Delta^2 + \varepsilon^2$,
hence positive unless $D^{+}$, $\Delta$ and $\varepsilon$ all vanish --- and there $\Lambda = 0$
anyway, so the guard can be placed on $\Delta$ rather than on the denominator. That is not a
cosmetic choice: writing it as a test on the denominator builds a relational around a
\code{Piecewise}, which the code generator folds into a ternary \emph{inside the condition of
another ternary}, and GCC then abandons the loop. The generated Venkatakrishnan kernel contained
zero vector instructions for exactly this reason until the guard was moved.

$\varepsilon^2$ is Venkatakrishnan's deactivation term, dimensional in the square of the solution
variable. Its purpose is to switch the limiter \emph{off} where the two-node variation is at the
level of round-off --- along flow-aligned edges in a channel core, where $u_i$ and $u_j$ agree to
round-off while the true face value legitimately differs from both transversally. Without it the
bound is decided by round-off there, the limiter clips a correction that was right, and refining
makes the error worse at $h^{-2}$. $\varepsilon^2 = 0$ recovers the preceding arm bit for bit,
which makes $K=0$ a genuine zero-severity control rather than an approximation of one.

\subsection{Scheme 3: Darwish--Moukalled}

van Leer's $\psi(r) = (r + |r|)/(1 + |r|)$ in the TVD form, with the limited increment
$\tfrac12\psi(r)(\phi_D - \phi_C)$. Substituting $r$ and clearing fractions
(\code{cvfem\_darwish\_moukalled\_inc}) gives
\begin{equation}
\Lambda \;=\; \frac{a\,|b| + |a|\,b}{2\,(|a| + |b|)},
\qquad a = 2\,\nabla\phi_C\cdot\dvec, \quad b = \phi_D - \phi_C ,
\label{eq:dm}
\end{equation}
which is one division rather than three and has no removable singularity: the denominator vanishes
only when $a$ and $b$ both do, and the numerator vanishes with it. Opposite signs give exactly
zero, which is $\psi(r<0) = 0$ --- the increment is discarded where the reconstruction and the
downwind difference disagree about the direction of the gradient. Note that this guard divides by
$2(|a|+|b|)$, an absolute value and not a \code{Piecewise}, which is why this kernel vectorised
when Venkatakrishnan's did not.

\subsection{Bound preservation, measured}

The schemes' claims are checked by instrumenting the limiter over a converged solve
($\sim$2M surface samples), counting how often the unlimited reconstruction leaves the two-node
interval, how often the limiter alters it, and how often the \emph{limited} value is still
outside:

\begin{center}
\begin{tabular}{@{}lrrr@{}}
\toprule
scheme & outside before & altered & outside after \\
\midrule
unlimited (control)        & 60.96\% & \phantom{0}0.00\% & 60.96\% \\
bounded-face clip          & 61.17\% & 61.17\% & 0.0000\% \\
Venkatakrishnan, $K=0$     & 61.11\% & 99.29\% & 0.0000\% \\
Darwish--Moukalled         & 61.13\% & 99.26\% & 0.0000\% \\
Venkatakrishnan, $K=5$     & 60.98\% & 65.84\% & 60.86\% \\
\bottomrule
\end{tabular}
\end{center}

The last row is the one to read carefully, and it is not a defect. With $\varepsilon^2 > 0$ the
limiter is deliberately inactive where the variation is at round-off, so it no longer enforces the
bound there --- "outside after" returns to the unlimited level. That is the trade $\varepsilon^2$
buys: bound preservation on the edges that matter, in exchange for not clipping the flow-aligned
ones where the two-node bound is meaningless. The two smooth limiters alter nearly every surface
($\sim$99\%) while the clip alters only those actually out of bounds, which is the difference
between a smooth rescaling and an active set.

\section{The micro-kernels}
\label{sec:mk}

This section describes how the components of Section~\ref{sec:components} are actually evaluated:
the staging, the loop structure, the compile-time dispatch, and what differs per scheme and per
layout. All counts below are properties of the written kernel and can be checked by reading it;
timings belong to the paper.

\subsection{Lane blocking and the packs}
\label{sec:mk:staging}

Every kernel processes a \emph{lane group} of elements at once. The width is
\code{CVFEM\_HEX8\_VEC\_SIZE = VEC\_BYTES / sizeof(scalar\_t)}, with \code{VEC\_BYTES} defaulting
to 128, so sixteen elements in double precision. Element data is staged into lane-major structs,
indexed \code{[node][lane]} --- or \code{[node][component][lane]} for the gradient packs --- so
that a fixed node index over varying lane is contiguous and every read in the surface loop is a
unit-stride vector load rather than a gather:

\begin{center}
\begin{tabular}{@{}lll@{}}
\toprule
pack & contents & read by \\
\midrule
\code{Hex8InputPack}    & $u_x,u_y,u_z,p$ per node            & every scheme \\
\code{Hex8ResidualPack} & the four accumulators per node      & every scheme \\
\code{Hex8RhieChowPack} & $\nabla p$ per node; $G_\scs$, $w_\scs$ per surface & Rhie--Chow arms \\
\code{Hex8UGradPack}    & $\nabla\uvec$ (nine components) and coordinates per node & higher-order arms \\
\bottomrule
\end{tabular}
\end{center}

\code{Hex8RhieChowPack} notably does \emph{not} carry node coordinates. It used to, purely to form
$\dvec$ by differencing; since the affine path takes $\dvec$ from the Jacobian
(Section~\ref{sec:geom:edges}) those three arrays have no reader, and removing them took 3\,KB of
stack per lane group and twenty-four indexed gathers per element out of the sweep. The
isoparametric kernels keep their own \code{Hex8CoordPack}.

\subsection{Loop structure}
\label{sec:mk:structure}

The residual's twelve surfaces are evaluated in \emph{one} lane loop
(\code{cvfem\_hex8\_conv\_all\_simd}), with each surface's contribution accumulated into
thirty-two lane-private scalars --- eight nodes times four rows --- and each stored once at the
end. The alternative, a vectorised lane loop per surface, writes a node's four accumulators three
times over, once per direction group, and measured slower.

The binding constraint on this loop is register capacity, and it is worth stating plainly because
it bounds what any rescheduling can achieve. One vector iteration holds thirty-two accumulators,
and each surface needs roughly two dozen more live values --- the area vector, both nodes'
velocity and pressure, and for the Rhie--Chow arms the pressure gradients and the two tabulated
surface scalars. That is on the order of fifty-six live vector values against the thirty-two
registers AArch64 provides, so the register allocator spills regardless of the order the surfaces
are visited in. Measured on the shipping binary, roughly $37\%$ of the convective sweep's
instructions were spill traffic against a $1952$-byte frame.

Three reschedulings were built and measured against this shape, and the pattern is instructive:
varying the lane width did nothing (the live set is per vector iteration, not across unrolled
iterations); fusing the surfaces as described above gained a few percent; and splitting the
momentum and continuity rows into separate passes --- which halves the accumulator count, exactly
the lever the spill count suggests --- \emph{lost} a third, because the per-surface inputs already
exceed what the accumulators leave free and the split only duplicated the mass-flux work. The
conclusion carried into Section~\ref{sec:geom:edges} is that the live set, not the schedule, is the
constraint, and that shrinking it means changing what a surface \emph{reads}.

\subsection{Compile-time dispatch}

Everything that is uniform over a sweep is a template parameter, not a runtime argument:

\begin{center}
\begin{tabular}{@{}ll@{}}
\toprule
flag & selects \\
\midrule
\code{RC}  & the Rhie--Chow term \\
\code{QG}  & the direction's pressure-gradient term in the Jacobian \\
\code{EPS} & the Harten band \\
\code{HO}  & the deferred correction \\
\code{LIM} & which limiter, $0..3$ \\
\code{NC}  & the field count in the reconstruction \\
\bottomrule
\end{tabular}
\end{center}

The reason is uniform across the file and was established by measurement rather than taste: a
sweep-uniform test left inside the vector body costs between $1.8\times$ and $2\times$, because it
tips the compiler's cost model out of vectorising the loop that contains it. A guard on the
Rhie--Chow coefficient cost $1.83\times$ this way; a runtime limiter select cost $2\times$. The
flags are dispatched once, outside the lane loop, by a macro that keeps the long argument list in
one place --- sixteen hand-written copies of it is where a transposed pack would hide.

\subsection{Per-scheme kernels}
\label{sec:mk:perscheme}

First-order upwinding runs the hand-written sweep with \code{HO=false}. The four higher-order
schemes are evaluated by \emph{generated} whole-element kernels, synthesised by
\code{python/synthesize\_cvfem\_hex8\_ns\_upwind\_sympy.py} and emitted into
\code{cvfem\_hex8\_ns\_upwind\_sympy\_kernels.hpp}: one per limiter, in two variants each with and
without Rhie--Chow, and all with a single \code{\#pragma omp simd} lane loop wrapping the entire
body:

\begin{center}
\begin{tabular}{@{}llrr@{}}
\toprule
scheme & kernel & lines, plain & lines, $+$RC \\
\midrule
unlimited          & \code{...\_defcor\_lim0\_simd} &  704 &  880 \\
bounded-face clip  & \code{...\_defcor\_lim1\_simd} &  798 &  963 \\
Venkatakrishnan    & \code{...\_defcor\_lim2\_simd} & 2290 & 2427 \\
Darwish--Moukalled & \code{...\_defcor\_lim3\_simd} & 1487 & 1611 \\
\bottomrule
\end{tabular}
\end{center}

The size ordering follows the limiter's algebra: the clip is a pair of comparisons, van Leer's
form has one division, and Venkatakrishnan's rational function with its deactivation term is three
times the body of either. Common subexpression elimination is applied across the whole element, so
the twelve surfaces share the centroid and geometry work that twelve separately compiled
per-surface functions could not.

These kernels contain no \code{if} statements at all --- the limiter logic is entirely ternaries,
which lower to selects and vectorise. That property is load-bearing and was not free: a ternary
whose \emph{condition} is itself a ternary does \emph{not} vectorise, because the compiler
if-converts the inner one into a PHI node and then cannot select over it. Venkatakrishnan's guard
had that shape (Section~\ref{sec:schemes}), and the kernel issued zero vector instructions until
the guard was moved onto the increment. The lesson generalises past this one kernel: in generated
code the property to check is not "is there a lane loop" --- all four always had one --- but
whether the compiler acted on it, and only a disassembly or a vectoriser report answers that.

\subsection{What the layout changes}

The two layouts differ in how an element's contribution reaches the global vector, not in the
arithmetic. The \emph{packed} layout gathers pack-locally, accumulates into pack-local
indices, and closes through a two-phase deterministic ghost reduction; the \emph{standard}
layout scatters with atomics. Consequences worth knowing when reading measurements:

\begin{itemize}
\item The packed layout is bitwise reproducible by construction. The standard layout is not ---
      its fingerprint varies run to run with the same binary, because the atomic order varies ---
      so a fingerprint comparison is only meaningful within the packed layout.
\item The packed layout spends a larger share of its time in the surface loop, so changes to the
      per-surface work show up there more strongly, and changes to per-element work (staging,
      gathers) show up more strongly on the standard layout.
\item Checksums are near-cancelling sums of order $10^{-14}$ built from terms of order one, so a
      last-bit difference reads as a large relative change. They are not a usable discriminator;
      the per-node cross-layout comparisons are.
\end{itemize}

\section{The selected kernel, per operation and layout}
\label{sec:sel}

This section is the dispatch, stated once. For each operation and each mesh layout it names the
function that runs, and the margin by which it beat the alternative that was built and measured
beside it. Ratios are between kernels measured against each other in \emph{one} job, with the job
id given, because the tree moved under these measurements: a ratio from one state is evidence
about a choice, an absolute throughput from the same job is not comparable with one from another.

\subsection{The dispatch}

\begin{center}
\footnotesize
\begin{tabular}{@{}llp{0.40\linewidth}@{}}
\toprule
operation & layout & kernel that runs \\
\midrule
residual, first order & packed
  & \code{apply\_residual\_packed} $\rightarrow$ \code{cvfem\_hex8\_ns\_upwind\_residual\_sumfact\_simd}
    $\rightarrow$ \code{cvfem\_hex8\_conv\_all\_simd} \\
residual, first order & standard
  & \code{apply\_residual\_atomic\_sumfact\_simd} $\rightarrow$ the same two \\
\midrule
residual, higher order & packed
  & \code{apply\_residual\_packed\_defcor}$(\texttt{sympy=true})$ $\rightarrow$ generated
    \code{cvfem\_hex8\_ns\_upwind\_sympy\_residual\_defcor\_rc\_lim$L$\_simd} \\
residual, higher order & standard
  & \code{apply\_residual\_atomic\_sumfact\_simd} $\rightarrow$ \emph{the same generated kernels},
    when $\texttt{venkat\_c} = 0$ \\
\midrule
Jacobian action & both
  & \code{cvfem\_hex8\_ns\_upwind\_jacobian\_action\_simd} $\rightarrow$
    \code{cvfem\_hex8\_conv\_all\_jv\_simd} $\rightarrow$ \code{cvfem\_hex8\_conv\_face\_jv\_simd} \\
Jacobian, partial assembly & both
  & \code{cvfem\_hex8\_ns\_upwind\_jacobian\_action\_pa\_simd} $\rightarrow$
    \code{cvfem\_hex8\_conv\_face\_jv\_pa\_simd} \\
\midrule
nodal gradients & packed & \code{cvfem\_hex8\_nodal\_grads\_packed\_nc<NC>} \\
nodal gradients & standard & \code{cvfem\_hex8\_nodal\_grads\_atomic\_nc<NC>} \\
\midrule
matrix assembly (BSR) & both & generated \code{cvfem\_hex8\_ns\_upwind\_sympy\_jacobian\_add\_bsr\_slots} \\
SpMV & --- & \code{bsr\_spmv\_static\_4x4} (\code{algebra/bsr/sfem\_BSR.hpp}) \\
\bottomrule
\end{tabular}
\end{center}

Two rows deserve emphasis. The higher-order residual runs the \emph{same} generated kernels on
both layouts --- there is no hand-vectorised second copy of that arithmetic, and the
\code{venkat\_c} condition is the only thing that diverts the standard layout to the hand-written
path. And the Jacobian names the same kernel for both layouts, because the layouts differ in how
the contribution reaches the global vector, not in the arithmetic.

\subsection{The surface-loop shape: opposite answers for residual and Jacobian}
\label{sec:sel:shape}

The twelve surfaces can be swept as one lane loop or as twelve, and the right choice is
\emph{not} the same for the two operators. This is the single most counter-intuitive selection in
the code and the reason both shapes exist.

\begin{center}
\footnotesize
\begin{tabular}{@{}lllrr@{}}
\toprule
operator & shape selected & alternative measured & packed & standard \\
\midrule
residual   & one fused lane loop, accumulate in registers
           & one vectorised loop per surface & $1.04$--$1.06$ & $1.04$ \\
residual   & \emph{(same)}
           & momentum/continuity in two passes & $0.66$ & $0.78$ \\
Jacobian   & one vectorised loop per surface
           & all twelve fused, into the pack & $0.78$ & $0.86$ \\
Jacobian   & \emph{(same)}
           & all twelve fused, into lane-private scalars & $0.88$ & $0.91$ \\
\bottomrule
\end{tabular}
\end{center}

\noindent
Residual rows from jobs \texttt{4953418} and \texttt{4954981}; Jacobian rows from job
\texttt{4952101}. Every one of these carried a control arm running untouched code, which read
within $0.4\%$ of unity --- without it the numbers would be unreadable, since node-to-node
variation on this machine is $4$--$5\%$.

The mechanism is register capacity, and it explains the sign change. Fusing lets a node's three
surface contributions be summed in a register and stored once, which is worth a few percent on
the residual. The Jacobian additionally carries the direction pack and the correction's
derivative, so fusing pushes its live set further past the thirty-two available registers than the
saved round-trips are worth. The row-split line is the sharpest evidence that the live set rather
than the schedule is binding: halving the accumulator count is exactly the lever the spill count
suggests, and it moved neither the spill frame ($1952 \rightarrow 1968$ bytes) nor the spill count
($1061 \rightarrow 1085$), while adding $20\%$ more instructions to recompute the mass flux.

\subsection{Generated against hand-written, per scheme}

For the higher-order schemes the generated whole-element kernels beat the hand-written lane
kernel on every arm (job \texttt{4954898}, packed, the \code{\_scalar} suffix naming the
hand-written path):

\begin{center}
\footnotesize
\begin{tabular}{@{}lll@{}}
\toprule
scheme & generated kernel & over hand-written \\
\midrule
unlimited, $+$Rhie--Chow & \code{...\_defcor\_rc\_lim0\_simd} & $1.48\times$ \\
bounded-face clip        & \code{...\_defcor\_lim1\_simd}     & $1.29\times$ \\
Venkatakrishnan          & \code{...\_defcor\_lim2\_simd}     & $1.27\times$ \\
Darwish--Moukalled       & \code{...\_defcor\_lim3\_simd}     & $1.20\times$ \\
\bottomrule
\end{tabular}
\end{center}

The margin is largest on the unlimited arm because that is where whole-element common
subexpression elimination has the most to share and the least limiter algebra to carry.

Venkatakrishnan's figure is the one to read with the history in mind. Until its guard was moved
off the denominator (Section~\ref{sec:schemes}) the generated kernel contained \emph{zero} vector
instructions, and fixing it was worth $1.21\times$ packed and $1.18\times$ standard on its own
(jobs \texttt{4953817} and \texttt{4954818}) --- which also moved the scheme from slower than
Darwish--Moukalled to level with it, the ordering its flop count predicts. A table of "best
kernels" assembled before that fix would have recorded the generated Venkatakrishnan kernel as the
best available while it was running scalar, which is a caution about what such a table does and
does not establish: it reports which of the \emph{built} alternatives won, not that the winner is
good.

\subsection{Where a rejected alternative is still in the tree}

Three selections above have alternatives that remain reachable, deliberately:

\begin{itemize}
\item The hand-written scalar higher-order path, reached by \code{--ho-scalar}. It is the oracle
      the generated kernels are verified against
      (\code{verify\_packed\_ho\_rc\_sympy\_vs\_scalar\_abs}, $\sim 3.5\times10^{-18}$), so it is
      a reference and not dead code.
\item The partially assembled Jacobian, selected by the solver rather than by performance: it
      trades a build for cheaper applies, which pays only because the Jacobian is applied many
      times at a fixed state.
\item The coloured layout (\code{cvfem\_hex8\_best\_colored.hpp}), which is neither the packed nor
      the atomic answer but a third scatter strategy kept for comparison.
\end{itemize}

\noindent
The reschedulings of Section~\ref{sec:sel:shape} that lost are \emph{not} in the tree. They are
recorded as comments with their numbers at the site that rejected them, because each one's static
argument --- a load count, a spill count, an accumulator count --- stays convincing enough that a
reader will otherwise rebuild them.

\section{The Jacobian action}
\label{sec:jac}

The solver applies $\partial R/\partial q$ to a direction $v$ without assembling it. Two variants
exist and they differ in whether the deferred correction and the Rhie--Chow reconstruction are
differentiated.

\subsection{Lagged and exact}

The \emph{lagged} action differentiates the first-order operator only: the correction
$\mathbf{h}_\scs$ contributes nothing, and the Rhie--Chow term is differentiated with the nodal
pressure gradient held fixed. This is the operator the assembled matrix can represent, because
including either term widens the pressure coupling past nearest neighbours and breaks the
one-ring sparsity.

The \emph{exact} action differentiates both. The costs are asymmetric and worth separating:

\begin{itemize}
\item Differentiating the Rhie--Chow correction, Eq.~\eqref{eq:corr}, requires the
      \emph{direction's} nodal pressure gradient, so a reconstruction pass runs per matvec. Its
      absence is not visible in the residual and does not make any single Newton step fail --- it
      caps Newton at a linear rate, and only where the continuity rows matter. On a
      backward-facing step, omitting it put the Jacobian $11.7\%$ away from a finite difference of
      the residual in the continuity rows while momentum stayed at $0.3\%$.
\item Differentiating the correction requires the direction's nodal \emph{velocity} gradient ---
      nine components against the three the pressure term needs --- plus the limiter's own
      derivative.
\end{itemize}

Both direction gradients are produced by one reconstruction sweep rather than two, since the pass
is parameterised by field count (\code{NC}).

\subsection{The coefficient's velocity sensitivity}

$G_\scs$ in Eq.~\eqref{eq:coeff} depends on the state through $|\bar{\uvec}|^2$, so the Jacobian
must carry $\partial G/\partial\uvec$. Writing $G = K/\sqrt{S}$ with
$K = \alpha|\Avec|^2/(\Avec\cdot\dvec)$ and $S$ the bracket of Eq.~\eqref{eq:coeff}, the
derivative along $v$ with face average $\bar{v}$ is
\begin{equation}
\delta G \;=\; -\,g\,\bigl(\bar{\uvec}\cdot\bar{v}\bigr),
\qquad
g \;=\; \frac{4\,u^2_{\text{scale}}}{h^2}\,\frac{G}{S} \;=\; G^3\,w ,
\end{equation}
with $w$ exactly Eq.~\eqref{eq:wdu}. The $G^3 w$ factorisation is the one used because $w$ is pure
geometry: it can be tabulated with the mesh, leaving the surface loop one multiply. The
Rhie--Chow mass flux $-G c$ therefore gains $g\,c\,(\bar{\uvec}\cdot\bar{v})$.

\subsection{Limiter derivatives and subgradients}
\label{sec:jac:subgrad}

The exact action needs $\partial\Lambda/\partial\Delta$, and two of the three limiters are not
differentiable everywhere. How that is handled is a correctness matter, not a detail.

\paragraph{The clip} is discontinuous in its derivative at both bounds. Naively differentiating
each branch makes the result depend on which side of the bound an ulp of round-off lands, and
since the two sides differ by a finite amount, an ulp flips an $O(1)$ quantity. In practice this
made the Jacobian action \emph{layout dependent}: the packed and standard arms disagreed at
$3\times10^{-3}$, because FMA contraction differs between them and tipped branches at ties. The
fix is a rounding band: within a tolerance scaled by the magnitudes involved, the derivative takes
the value that does not depend on which side the comparison fell.

\paragraph{Darwish--Moukalled} has the same defect at $a=0$ and $b=0$, where the absolute values
have no derivative. The two one-sided limits differ by $2\,\delta a\,b/\mathrm{den}$ and
$2\,a\,\delta b/\mathrm{den}$ respectively. Inside a rounding band the implementation takes the
\emph{minimum-norm subgradient}, which is zero --- the symmetric element of the subdifferential of
$|\cdot|$ at zero. Neither side is preferred, and it is the only choice independent of where an
ulp lands.

\paragraph{Venkatakrishnan} branches too, on $\Delta \ge 0$, but needs no band: at $\Delta = 0$ its
$\psi$ is exactly $1$ and every term carrying the branch is multiplied by $\Delta$, so the
discontinuity cancels and the derivative is continuous there. This was checked rather than assumed.

\subsection{Partial assembly}

A third variant precomputes the surface quantities that do not depend on the direction and stores
them, so the apply reads a compressed tangent instead of the state
(\code{cvfem\_hex8\_build\_pa\_tangent}). Per surface it stores the upwind split, and for the
Rhie--Chow arms the combined velocity-sensitivity factor $g\,c\,\bar{\uvec}$ as three components.
The store is sixty scalars per element.

This is where the economics of precomputation invert relative to the residual. The Jacobian is
applied many times at a \emph{fixed} state, so the build amortises over the Krylov iterations; the
residual is evaluated once per state, so precomputing for it would mean rebuilding every call.
That asymmetry is why the Jacobian reads a tabulated $G_\scs$ while the residual recomputes it
inline, and why the same change --- taking $\dvec$ from the Jacobian --- helped the residual
substantially and the Jacobian not at all until $w$ was tabulated too.

\subsection{Verification}

Three oracles, each catching a different class of error:

\begin{itemize}
\item \emph{Cross-layout}: the packed and standard actions compared node by node, which catches
      staging and scatter errors. Reads $\sim10^{-16}$.
\item \emph{Finite difference of the residual}: catches a missing or wrong term in the
      linearisation. The exact arm is gated; the lagged arm is gated \emph{from below}, because a
      lagged action that agreed with the finite difference would mean the correction never reached
      the residual and the comparison is vacuous.
\item \emph{Partially assembled against direct}: catches disagreement between a producer of
      precomputed data and its consumer. This is the only one of the three that runs on a
      non-affine mesh, and consequently the only one that can see a geometry inconsistency ---
      Section~\ref{sec:geom:edges} records it catching one at $1.4\times10^{-5}$ that the other
      two could not.
\end{itemize}

A verification must include whichever passes the timed apply included. Comparing an action that
ran the boundary closure against a reference that did not reports the closure as a staging error,
which has happened and reads as an $O(1)$ failure.

\section{The nodal gradient reconstruction}
\label{sec:recon}

Rhie--Chow needs a nodal pressure gradient, and the exact higher-order Jacobian needs a nodal
velocity gradient --- the state's for the residual's correction, the direction's for its
derivative. One pass produces all of them
(\code{cvfem\_hex8\_nodal\_grads\_packed\_nc} and its atomic twin in
\code{src/hex8/cvfem\_hex8\_pack\_helpers.hpp}).

\subsection{What it computes}

For each element the reference gradient of the field is contracted from the nodal values and
pushed forward by $\adj\Jmat/\det\Jmat$, giving one element-constant gradient per field
component; that value is then distributed to all eight of the element's nodes with a weight, and
the nodal gradient is the weighted sum over the elements meeting at the node. The route ---
reference gradient, then geometry --- is the same two-stage shape every other gradient in the code
uses, deliberately, so there is one implementation of the mapping.

Three properties shape the kernel:

\begin{itemize}
\item \emph{The element gradient does not depend on the node.} The same $n_c$-vector is
      broadcast to all eight nodes, so the pass is scatter-dominated: at three fields that is
      seventy-two scalar read-modify-writes per element against roughly sixty-six flops of real
      work.
\item \emph{It is a reduction over elements.} Hence the same layout question as the operator:
      pack-local accumulation with a deterministic ghost closure, or atomics.
\item \emph{The field count is a template parameter} (\code{NC}). A runtime trip count in the
      scatter cost measurably, so the single code path is instantiated per field count rather than
      parameterised at runtime.
\end{itemize}

\subsection{One path, two field counts}

The pass is used at $n_c = 1$ (the pressure gradient, in every Rhie--Chow residual) and $n_c = 3$
(the velocity gradient, per matvec in the exact higher-order Jacobian). These were once two
separate implementations, and unifying them removed more than duplication: the multi-field path
could not read strided input, which was the only reason the exact action de-interleaved the
direction's three velocity components into scratch on every matvec. Output is described by a
per-component pointer and stride, input by a stride, so the single kernel reads the interleaved
direction in place.

The unification measured $+10\%$ on the packed layout and $+27\%$ on the standard one at both field
counts. Performance neutrality at $n_c = 1$ was the gate that had to be cleared before anything
else, since that path runs in every Rhie--Chow residual.

\subsection{Rejected optimisations, and why they are recorded}

Three further changes to this pass were built, measured and reverted. They are recorded in the
source with their numbers, because each is the kind of idea that looks obviously right on a static
count and is wrong in fact:

\begin{itemize}
\item \emph{Templating the geometry kind} to hoist the isoparametric branch out of the element
      loop. Lost; and the first diagnosis of why was also wrong, blaming an OpenMP capture that
      turned out to be innocent.
\item \emph{Writing owned rows straight out}, skipping the pack-local buffer for the contiguous
      owned range. The static argument is sound --- it removes a read-and-write pass over the
      owned majority --- and it measured slower.
\item \emph{Lane-blocking the element loop}, which every operator sweep beside it already does.
      Lost at this pass's arithmetic-to-scatter ratio.
\end{itemize}

The reason to keep the numbers rather than the conclusions is that all three arguments remain
superficially convincing, and a reader who re-derives them from a load count will rebuild them.
The same applies to the surface-loop reschedulings of Section~\ref{sec:mk:structure}.

\section{Reading the source from here}
\label{sec:map}

\begin{center}
\begin{tabular}{@{}ll@{}}
\toprule
component & file \\
\midrule
surfaces, geometry, fluxes, packs & \code{src/upwind/cvfem\_hex8\_ns\_upwind\_kernels.hpp} \\
generated higher-order kernels    & \code{src/upwind/cvfem\_hex8\_ns\_upwind\_sympy\_kernels.hpp} \\
their generator                   & \code{python/synthesize\_cvfem\_hex8\_ns\_upwind\_sympy.py} \\
limiters and their derivatives    & \code{src/venkata/cvfem\_venkata\_limiter.hpp} \\
coefficient table, PA tangent, reconstruction & \code{src/hex8/cvfem\_hex8\_pack\_helpers.hpp} \\
packed / standard / coloured sweeps & \code{src/best/cvfem\_hex8\_best\_*.hpp} \\
\bottomrule
\end{tabular}
\end{center}

\noindent
Two conventions that make the source easier to read. Anything uniform over a sweep is a template
parameter, so a function's behaviour is determined by its instantiation and not by its arguments;
and the hand-written scalar kernels are kept as references for the vectorised and generated ones
to be measured against, not as dead code --- they are what the oracles of
Section~\ref{sec:jac} compare against.


\end{document}
